\section{Discussion}

\subsection{Why the Proxy Failed}

Canonical HumanEval solutions are carefully crafted reference implementations. They pass all tests by design and follow Python conventions. LLM-generated code---especially early attempts in a repair loop---contains more issues: undefined variables, type mismatches, incomplete logic.

Using canonical solutions as an LLM proxy conflates the reference answer with the model's attempts. The 9.15\% warning rate reflects human code quality, not LLM generation behavior.

\subsection{What the Baseline Tells Us}

Despite the proxy limitation, our results provide value:

\begin{enumerate}
    \item \textbf{Floor established:} Canonical solutions have $\sim$9\% warning rate. LLM code should exceed this, providing room for the static analysis intervention to act.

    \item \textbf{Pipeline validated:} The pylint analysis infrastructure works. Components are reusable for future work with LLM API access.

    \item \textbf{Warning types identified:} The 9 warning categories are actionable for code repair. They explain \emph{why} code is problematic.
\end{enumerate}

\subsection{Gate Design Implications}

The MUST\_WORK gate served its purpose. Had we skipped the existence check and proceeded to full pass@k evaluation, we would have:
\begin{itemize}
    \item Wasted compute on the wrong code source
    \item Produced results that don't test the hypothesis
    \item Potentially drawn incorrect conclusions
\end{itemize}

Gate-based validation catches methodology issues early.

\subsection{Limitations}

\textbf{Primary:} No LLM-generated code tested. API access unavailable during this study.

\textbf{Secondary:}
\begin{itemize}
    \item Single dataset (HumanEval only, not MBPP)
    \item Single static analyzer (pylint, not mypy)
    \item Single severity filter configuration
\end{itemize}

\textbf{Framing:} This is a methodology study establishing baseline and validating infrastructure---not a full hypothesis test.

\subsection{Path Forward}

To complete the hypothesis evaluation:

\begin{enumerate}
    \item \textbf{Re-run h-e1 with LLM API:} Generate code using GPT-4 or Claude, measure warning rates
    \item \textbf{If h-e1 passes:} Proceed to h-m1 (LLM interprets warnings) and h-m2 (non-overlapping errors)
    \item \textbf{Full evaluation:} Compare pass@k between static+execution vs.\ execution-only conditions
\end{enumerate}

The pipeline, baseline, and methodology are ready. The missing ingredient is LLM API access.
